%\documentclass[a4paper,12pt]{article}
%\documentclass[a4paper,landscape,11pt]{report}
\documentclass[a4paper,11pt]{article}
%
\usepackage{amsmath}
\usepackage{amssymb}
\usepackage{array}
\usepackage[french]{babel}
\usepackage{calc}
%\usepackage{caption}
\usepackage[hang]{caption}
\usepackage{cite}
\usepackage{color}
\usepackage{enumerate}
\usepackage{epic}
\usepackage{eepic}
\usepackage{eepicemu}
\usepackage{epsfig}
\usepackage{fancyhdr}
\usepackage{float}
\usepackage[T1]{fontenc}
%\usepackage{fourier}
\usepackage[dvips]{geometry}
\usepackage{graphicx}
\usepackage{import} % To use \subimport and import stuff 
\usepackage[utf8]{inputenc}
%\usepackage{kpfonts}
\usepackage{lastpage}
\usepackage{lmodern}
\usepackage{multicol}
\usepackage{rotating}
\usepackage{setspace}
%\usepackage{subfigure} %% deprecated ?
\usepackage{subcaption}
\usepackage{verbatim}
\usepackage{dsfont}
\usepackage{gensymb}







\geometry{top=1.5cm, bottom=1.5cm, left=1.5cm, right=1.5cm}

%\newcommand\textsubscript[1]{\ensuremath{{}_{\text{#1}}}}
\newcounter{toto}
\newcommand{\pn}{\par\noindent}
\newcommand{\vs}{\vskip 0.5 \baselineskip}
\newcommand{\la}{\langle}
\newcommand{\ra}{\rangle}
%
\newcommand{\HRule}{\rule{\linewidth}{0.5mm}}
\newcommand{\bm}{\boldmath}
\newcommand{\um}{\unboldmath}
\newcommand{\hm}{\fontfamily{phv}\fontseries{m}\selectfont}
\newcommand{\hb}{\fontfamily{phv}\fontseries{b}\selectfont}
\newcommand{\cf}[1]{{\tt (cf. : #1)} }
%
\def\div{\mathop{\rm div}\nolimits}
\def\grad{\mathop{\overrightarrow{\rm grad}}\nolimits}
\def\rot{\mathop{\overrightarrow{\rm rot}}\nolimits}
\def\u{\mathop{\vec u}\nolimits}
%
%
%
%%%%%%%%%%%%%%%%%%%%%%%%%%%%%%%%%%%%%%%%%%%%%%%%%%%%%%%%%%%%%%%%%%%%%%%%%%%%%
%
%
%\renewcommand{\chaptername}{Lecture}
\addto\captionsfrench{\renewcommand{\chaptername}{Partie}}

\linespread{1.1}
\begin{document}
%
\thispagestyle{empty}
\renewcommand{\headrulewidth}{0.0pt}
%

\setlength{\columnsep}{30pt}
\setlength{\columnseprule}{1pt}
%\begin{multicols}{2}
\noindent
%\begin{minipage}[t]{0.35\textwidth}
\begin{minipage}[t]{0.35\textwidth}
~\\[-2ex]
% \includegraphics[width=1.8cm]{/home/spetit/Administration/Logos/logo-sciences-couleur-texte_new}
\includegraphics[scale=0.7]{unicaen_logo_rvb_noir_V1}

\end{minipage}
%\begin{minipage}[t]{0.5\textwidth}
\begin{minipage}[t]{0.50\textwidth}
\flushleft{
\footnotesize{UFR des sciences} \\
%\tiny{~}\\
\footnotesize{Licence PCI -- seconde année} \\
%\tiny{~}\\
\footnotesize{Initiation aux méthodes numériques}
}
\end{minipage}
%\hfill
\begin{minipage}[t]{0.12\textwidth}
\flushleft
{
\footnotesize{2017-2018}}
\end{minipage}

\vspace{0.5cm}

\centerline{\sc\Large Travaux pratiques de programmation}
{\let\clearpage\relax \section*{Introduction}}
Les travaux pratiques qui suivent vise à vous aider à mieux appréhender 
les méthodes de raisonnements et les outils de programmation. 
Ceux-ci étant largement indispensable à la pratique des sciences à
haut niveau, dans laquelle il est souvent indispensable
de créer ses propres outils. 

Pour des raisons d'accessibilités la programmation sera faite en Python ;
ce langage étant pratiqué au lycée. Le langage lui-même ne sera pas enseigné,
vous serez amenés à en développer une connaissance empirique par vous mêmes. 
Un fort niveau d'engagement personnel est donc nécessaire en séance. 

Chaque TP aborde des notions simples de mathématique à aborder avec 
les éléments les plus simples de Python. Le tout est conçu pour vous
permettre de vous concentrer sur l'aspect essentiel de ces TP :
l'algorithmique.

% 1 séances 2 exercices
%\newcommand{cf}[1]{({\sc cf. :}{\tt #1.})}
\subimport{./tp_derivation/exercices/}{enonce.tex}
%\section{Derivation}
Il arrive que la dérivation analytique d'une fonction ne soit pas facilement
calculabe. Ce TP survol des méthodes élémentaires de dérivation numérique.
\subsection{Précision}
Cet exercice vise à étudier les erreurs commises lors de l'emploi
d'un méthode de dérivation numérique.
\begin{enumerate}
\item Écrire un fonction capable de calculer la dérivée 
d'une fonction (de $\mathbb{R}  \longrightarrow  \mathbb{R}$)  en 
un point $x$ quelconque en utilisant la formule : 
$$f'(x)= \frac{f(x+h)-f(x-h)}{2h}\,;$$
où $h$ est un paramètre qui pourra être choisi lors de l'appel de cette fonction
\cf{function.py}.
\item Utiliser cette fonction pour calculer $\frac{d\cos(x)}{dx}$ 
sur cent points répartis homogènements sur $[0\,;\,\,2\pi]$
\cf{mesh.py}.
Ensuite tracer la fonction $\cos(x)$ et les points calculés sur
un graphique
\cf{plot.py}.
\item Sur les cent points, calculer la moyenne de la valeur
 absolue de la différence entre dérivée numérique et fonction analytique
\cf{loop.py}.
\item Avec une échelle log/log, tracer l'évolution de cette différence 
en fonction de $h$. Les valeurs choisies pour $h$ suivront une suite 
géométrique de raison $2^\frac{1}{2}$ et de valeur initiale $10^{-6}$ \cf{loop.py}. 
On utilisera les 25 premières valeurs.
Qu'observe-t-on ? 
Comment choisir la meilleure valeur de $h$ pour dériver numériquement ? 
\item Tracer la dérivée numérique (pour une valeur de $h$ choisie arbitrairement)
 et la dérivée 
analytique sur un même graphique. Peut-on observer facilement la différence ? 
\item Tracer la différence ente dérivée numérique et calcul analytique. 
\item Comparer la forme de la fonction obtenue avec 
$-\alpha \sin(x)$ en choisissant de manière appropriée $\alpha$.
Expliquer ce qui est observé à partir du développement limité de $\cos(x)$ 
autour de $0$.. 
\end{enumerate}

\subsection{Algorithmes de dérivation}
De nombreux algorithmes de dérivation numérique existent. Cette section vise à
en programmer quelques-uns et à les comparer.
\begin{enumerate}
\item Programmer les fonctions de dérivation utilisant les formules : 
$$f'(x)= \frac{-f(x+2h)+8f(x+h)-8f(x-h)+f(x-2h)}{12h}\,;$$ 
$$f'(x)= \frac{f(x+3h)-9f(x+2h)+45f(x+h)-45f(x-h)+9f(x-2h)-f(x-3h)}{60h}.$$
Les coefficients de ces fonctions sont calculés de manière à annuler de plus en
plus de termes du développement limité de la fonction $f$.
\item Des formules non-symétriques existent également. En voici quelques-unes :
 $$f'(x) = \frac{f(x+h)-f(x)}{h}\,;$$
$$f'(x)=\frac{-3f(x)+4f(x+h) -f(x+2h)}{2h}\,;$$
$$f'(x)=\frac{-15f(x)+16f(x+h) -f(x+4h)}{12h}\,;$$
 Programmer les fonctions de dérivation correspondantes. Ces fonctions peuvent être utilisées avec $h>0$ ou $h<0$.
\item À nombre d'évaluation de la fonction $f$ égale, comparer la précision moyenne de
chacune de ces méthodes sur un graphique (pour des valeurs de $h$ comprises entre $0,001$ et $1$).
\end{enumerate}



\subsection{Ordres supérieurs}
Il est possible d'obtenir des dérivées d'ordre supérieur par la même méthode (seconde, troisième\ldots).
\begin{enumerate}
\item Programmer une fonction calculant la dérivée à l'ordre 2 d'une fonction en 
utilisant l'algorithme $f'(x)= \frac{f(x+h)-f(x-h)}{2h}\,;$ pour obtenir 
des dérivées première, puis calculant la dérivée seconde à partir de deux points ainsi obtenus. 
\item Programmer une fonction calculant la dérivée seconde à l'aide de la formule :
$$f'(x)= \frac{f(x+h)+f(x-h)-2f(x)}{h^2}.$$
Comparer l'évolution de l'erreur de ces deux méthodes en fonction de $h$. Puis réaliser
une comparaison \og{} sur papier \fg{} des deux méthodes. 
\end{enumerate}

\subsection{Dérivées de série de points régulière}
Ces méthodes de dérivation numérique peuvent être utiles pour dériver des fonctions 
dont la forme analytique est complexe et pour lesquelles on dispose juste 
d'une série de points. À titre d'exercice nous allons générer une telle série 
de points dans un fichier, puis calculer la dérivée.

\begin{enumerate}
\item Écrire un programme générant un fichier contenant deux colonnes \cf{file.py}. Dans la première
on trouvera les valeurs de $x$ et dans la seconde celles de $f(x)$ ; où $f$ sera une fonction de votre choix. Générer 100 points sur un intervalle $[a\,;\,b]$ de votre choix. 
\item Écrire un programme capable de relire ce fichier.
\item Combien de valeurs de la dérivée peut-on calculer avec une formule de dérivation centrée ? 
\item Écrire une fonction capable de calculer la dérivée de la fonction échantillonnée dans le fichier pour chacun des points.
\end{enumerate}

\subsection{Dérivées de séries de points irrégulière}
Reproduire le même travail, mais cette fois avec une série de points distribués de manière irrégulière sur l'intervalle $[a\,;\,b]$. Il faudra donc calculer à chaque fois de nouveaux coefficients pour la dérivation. 

%\section{Derivation}
Il arrive que la dérivation analytique d'une fonction ne soit pas facilement
calculabe. Ce TP survol des méthodes élémentaires de dérivation numérique.
\subsection{Précision}
Cet exercice vise à étudier les erreurs commises lors de l'emploi
d'un méthode de dérivation numérique.
\begin{enumerate}
\item Écrire un fonction capable de calculer la dérivée 
d'une fonction (de $\mathbb{R}  \longrightarrow  \mathbb{R}$)  en 
un point $x$ quelconque en utilisant la formule : 
$$f'(x)= \frac{f(x+h)-f(x-h)}{2h}\,;$$
où $h$ est un paramètre qui pourra être choisi lors de l'appel de cette fonction
\cf{function.py}.
\item Utiliser cette fonction pour calculer $\frac{d\cos(x)}{dx}$ 
sur cent points répartis homogènements sur $[0\,;\,\,2\pi]$
\cf{mesh.py}.
Ensuite tracer la fonction $\cos(x)$ et les points calculés sur
un graphique
\cf{plot.py}.
\item Sur les cent points, calculer la moyenne de la valeur
 absolue de la différence entre dérivée numérique et fonction analytique
\cf{loop.py}.
\item Avec une échelle log/log, tracer l'évolution de cette différence 
en fonction de $h$. Les valeurs choisies pour $h$ suivront une suite 
géométrique de raison $2^\frac{1}{2}$ et de valeur initiale $10^{-6}$ \cf{loop.py}. 
On utilisera les 25 premières valeurs.
Qu'observe-t-on ? 
Comment choisir la meilleure valeur de $h$ pour dériver numériquement ? 
\item Tracer la dérivée numérique (pour une valeur de $h$ choisie arbitrairement)
 et la dérivée 
analytique sur un même graphique. Peut-on observer facilement la différence ? 
\item Tracer la différence ente dérivée numérique et calcul analytique. 
\item Comparer la forme de la fonction obtenue avec 
$-\alpha \sin(x)$ en choisissant de manière appropriée $\alpha$.
Expliquer ce qui est observé à partir du développement limité de $\cos(x)$ 
autour de $0$.. 
\end{enumerate}

\subsection{Algorithmes de dérivation}
De nombreux algorithmes de dérivation numérique existent. Cette section vise à
en programmer quelques-uns et à les comparer.
\begin{enumerate}
\item Programmer les fonctions de dérivation utilisant les formules : 
$$f'(x)= \frac{-f(x+2h)+8f(x+h)-8f(x-h)+f(x-2h)}{12h}\,;$$ 
$$f'(x)= \frac{f(x+3h)-9f(x+2h)+45f(x+h)-45f(x-h)+9f(x-2h)-f(x-3h)}{60h}.$$
Les coefficients de ces fonctions sont calculés de manière à annuler de plus en
plus de termes du développement limité de la fonction $f$.
\item Des formules non-symétriques existent également. En voici quelques-unes :
 $$f'(x) = \frac{f(x+h)-f(x)}{h}\,;$$
$$f'(x)=\frac{-3f(x)+4f(x+h) -f(x+2h)}{2h}\,;$$
$$f'(x)=\frac{-15f(x)+16f(x+h) -f(x+4h)}{12h}\,;$$
 Programmer les fonctions de dérivation correspondantes. Ces fonctions peuvent être utilisées avec $h>0$ ou $h<0$.
\item À nombre d'évaluation de la fonction $f$ égale, comparer la précision moyenne de
chacune de ces méthodes sur un graphique (pour des valeurs de $h$ comprises entre $0,001$ et $1$).
\end{enumerate}



\subsection{Ordres supérieurs}
Il est possible d'obtenir des dérivées d'ordre supérieur par la même méthode (seconde, troisième\ldots).
\begin{enumerate}
\item Programmer une fonction calculant la dérivée à l'ordre 2 d'une fonction en 
utilisant l'algorithme $f'(x)= \frac{f(x+h)-f(x-h)}{2h}\,;$ pour obtenir 
des dérivées première, puis calculant la dérivée seconde à partir de deux points ainsi obtenus. 
\item Programmer une fonction calculant la dérivée seconde à l'aide de la formule :
$$f'(x)= \frac{f(x+h)+f(x-h)-2f(x)}{h^2}.$$
Comparer l'évolution de l'erreur de ces deux méthodes en fonction de $h$. Puis réaliser
une comparaison \og{} sur papier \fg{} des deux méthodes. 
\end{enumerate}

\subsection{Dérivées de série de points régulière}
Ces méthodes de dérivation numérique peuvent être utiles pour dériver des fonctions 
dont la forme analytique est complexe et pour lesquelles on dispose juste 
d'une série de points. À titre d'exercice nous allons générer une telle série 
de points dans un fichier, puis calculer la dérivée.

\begin{enumerate}
\item Écrire un programme générant un fichier contenant deux colonnes \cf{file.py}. Dans la première
on trouvera les valeurs de $x$ et dans la seconde celles de $f(x)$ ; où $f$ sera une fonction de votre choix. Générer 100 points sur un intervalle $[a\,;\,b]$ de votre choix. 
\item Écrire un programme capable de relire ce fichier.
\item Combien de valeurs de la dérivée peut-on calculer avec une formule de dérivation centrée ? 
\item Écrire une fonction capable de calculer la dérivée de la fonction échantillonnée dans le fichier pour chacun des points.
\end{enumerate}

\subsection{Dérivées de séries de points irrégulière}
Reproduire le même travail, mais cette fois avec une série de points distribués de manière irrégulière sur l'intervalle $[a\,;\,b]$. Il faudra donc calculer à chaque fois de nouveaux coefficients pour la dérivation. 





\end{document}


